% Generated by build_cv.py from the website's YAML data.
{\fontsize{26}{29}\selectfont Kowndinya Boyalakuntla}\par\vspace{5pt}
{\sffamily Ph.D. candidate in Computer Science \enspace\textbar\enspace Rutgers University}\par\vspace{5pt}
{\small \href{mailto:kb1204@scarletmail.rutgers.edu}{kb1204@scarletmail.rutgers.edu}\enspace\textbar\enspace \href{https://kowndinya2000.github.io}{kowndinya2000.github.io}\enspace\textbar\enspace \href{https://scholar.google.com/citations?user=tv9b9Y4AAAAJ}{Google Scholar}}\par
{\sffamily\scriptsize Updated October 2026}\par
\cvsection{Introduction}
I am a Ph.D. candidate in Computer Science at Rutgers University with over three years of research in robot learning and planning. I have published at CoRL and IROS, built robotics benchmarks and evaluation pipelines, and contributed to diffusion-based 6D object placement and language-guided rearrangement. My experience spans 3D scene understanding and sim-to-real manipulation. My current research focuses on learned world models and Vision-Language-Action (VLA) models for humanoid loco-manipulation.
\cvsection{Education}
\educationentry{Rutgers University}{Sep 2024 -- present}{Ph.D. in Computer Science; GPA: 3.85/4.0}{Advisor: Prof. Abdeslam Boularias}
\educationentry{Rutgers University}{Sep 2022 -- Oct 2024}{M.S. in Computer Science; GPA: 3.8/4.0}{Advisor: Prof. Abdeslam Boularias\par {\small Thesis: Context-Aware Entity Grounding with Open-Vocabulary 3D Scene Graphs.}}
\educationentry{Indian Institute of Technology Tirupati}{Jul 2017 -- Jul 2021}{B.Tech. in Computer Science and Engineering; GPA: 8.42/10.0}{Advisors: Prof. Sridhar Chimalakonda, Prof. Kalidas Yeturu}
\cvsection{Research experience}
\cventry{Rutgers University}{Jan 2023 -- present}{Graduate Student Researcher · New Brunswick, NJ}{Conducted research in robot learning for 3D scene understanding, object rearrangement, and placement.}
\cventry{RISHA Lab, IIT Tirupati}{Dec 2019 -- May 2022}{Undergraduate Student Researcher · Tirupati, India}{Developed tools for GitHub repository evaluation and energy-aware mobile application generation.}
\cventry{University of Waterloo}{Jul 2020 -- Sep 2020}{Visiting Scholar · Waterloo, Canada}{Improved RepoReaper's recall and reduced false positives in GitHub repository analysis without GHTorrent.}
\cvsection{Publications}
\cvsubsection{First-Author}
\begin{papers}
\paper{https://arxiv.org/abs/2609.39751}{PL-MPC: Beyond Policy Alignment: Closing the Planning--Learning Loop for Robot Control with Learned World Models}{\textbf{Kowndinya Boyalakuntla}, Yuhan Liu, Abdeslam Boularias}{Manuscript, 2026. Under review.}{\resource{https://arxiv.org/abs/2609.39751}{Paper}\enspace \resource{https://pl-mpc-humanoid.github.io/}{Project}\enspace \resource{https://github.com/Kowndinya2000/pl-mpc}{Code}}{\begin{itemize}[leftmargin=12pt,label=\textbullet,itemsep=1pt,parsep=0pt,topsep=2pt]\item Combined hybrid multi-step value targets, uncertainty-aware MPC, and return-weighted policy distillation to improve planning and learning with a shared world model.
\item Improved mean return over TD-M(PC)\textsuperscript{2} from 98 to 387 on HumanoidBench balance-hard and from 199 to 466 on hurdle; demonstrated zero-shot transfer to KUKA wrench--nut alignment.\end{itemize}}
\paper{https://arxiv.org/abs/2609.38857}{Plan-Conditioned Imitation for Robust Object Retrieval under Self-Occlusion in Dense Clutter}{\textbf{Kowndinya Boyalakuntla}, Ajinkya Pawar, Abdeslam Boularias, Jingjin Yu}{Manuscript, 2026. Under review.}{\resource{https://arxiv.org/abs/2609.38857}{Paper}\enspace \resource{https://trace-retrieval.github.io/}{Project}\enspace \resource{https://github.com/Kowndinya2000/trace-code}{Code}\enspace \resource{https://huggingface.co/datasets/Kowndi/trace}{Data}}{\begin{itemize}[leftmargin=12pt,label=\textbullet,itemsep=1pt,parsep=0pt,topsep=2pt]\item Trained a recurrent policy with behavior cloning and DAgger to combine a fixed teacher plan with partial observations for object retrieval under self-occlusion.
\item Achieved 90.7\% success on 511 simulation scenes and 90\% on UR5e; mean hardware task time was 67.3 s versus 192.7 s for the privileged teacher, which achieved 95\% success.\end{itemize}}
\paper{https://ieeexplore.ieee.org/abstract/document/11245879}{KARL: Kalman-Filter Assisted Reinforcement Learner for Dynamic Object Tracking and Grasping}{\textbf{Kowndinya Boyalakuntla}, Abdeslam Boularias, Jingjin Yu}{IROS, 2025}{\resource{https://ieeexplore.ieee.org/abstract/document/11245879}{Paper}\enspace \resource{https://github.com/arc-l/karl}{Code}}{\begin{itemize}[leftmargin=12pt,label=\textbullet,itemsep=1pt,parsep=0pt,topsep=2pt]\item Integrated extended Kalman filtering for tracking recovery, a six-stage reinforcement learning curriculum, and grasp retries for dynamic grasping with an eye-in-hand camera.
\item Achieved 93.3\% versus 60.0\% grasp success in regular scenes and 80.0\% versus 53.3\% in obstacle scenes against EARL on UR5e, with 15 trials per condition.\end{itemize}}
\end{papers}
\cvsubsection{Co-Author}
\begin{papers}
\paper{https://kowndinya2000.github.io/pdfs/DAP\_IROS24.pdf}{DAP: Diffusion-based Affordance Prediction for Multi-modality Storage}{Haonan Chang, \textbf{Kowndinya Boyalakuntla}, Yuhan Liu, Xinyu Zhang, Liam Schramm, Abdeslam Boularias}{IROS, 2024}{\resource{https://kowndinya2000.github.io/pdfs/DAP\_IROS24.pdf}{Paper}\enspace \resource{https://github.com/changhaonan/DPS}{Code}}{\begin{itemize}[leftmargin=12pt,label=\textbullet,itemsep=1pt,parsep=0pt,topsep=2pt]\item Collected data and evaluated a pipeline combining diffusion-based affordance sampling with point-cloud correspondence for 6D object placement.
\item Achieved 98\% book-shelving and 94\% can-stacking success in simulation, versus RPDiff's 94\% and 85\%; demonstrated real-robot fruit storage after training on 80 demonstrations.\end{itemize}}
\paper{https://arxiv.org/abs/2309.15821}{LGMCTS: Language-Guided Monte-Carlo Tree Search for Executable Semantic Object Rearrangement}{Haonan Chang, Kai Gao, \textbf{Kowndinya Boyalakuntla}, Alex Lee, Baichuan Huang, Harish Udhaya Kumar, Jingjin Yu, Abdeslam Boularias}{IROS, 2024}{\resource{https://arxiv.org/abs/2309.15821}{Paper}\enspace \resource{https://lgmcts.github.io}{Project}\enspace \resource{https://github.com/changhaonan/LGMCTS}{Code}}{\begin{itemize}[leftmargin=12pt,label=\textbullet,itemsep=1pt,parsep=0pt,topsep=2pt]\item Co-developed ELGR-Bench and implemented object selection, language parsing, and evaluation for a planner combining language-derived geometric priors with Monte Carlo tree search.
\item Achieved 79.2\% executable-plan success on ELGR-Bench versus 45.2\% for LLMs as Few-Shot Planners and 74.1\% for separate goal-pose and motion planning; demonstrated rearrangement on UR5e.\end{itemize}}
\paper{https://arxiv.org/abs/2309.15940}{Context-Aware Entity Grounding with Open-Vocabulary 3D Scene Graphs}{Haonan Chang, \textbf{Kowndinya Boyalakuntla}, Shiyang Lu, Siwei Cai, Eric Pu Jing, Shreesh Keskar, Shijie Geng, Adeeb Abbas, Lifeng Zhou, Kostas Bekris, Abdeslam Boularias}{CoRL, 2023}{\resource{https://arxiv.org/abs/2309.15940}{Paper}\enspace \resource{https://ovsg-l.github.io}{Project}\enspace \resource{https://github.com/changhaonan/OVSG}{Code}}{\begin{itemize}[leftmargin=12pt,label=\textbullet,itemsep=1pt,parsep=0pt,topsep=2pt]\item Built the DOVE-G benchmark (8 scenes, 4,000 queries) and evaluation pipeline for grounding language in open-vocabulary 3D scene graphs; reconstructed scenes and integrated ConceptFusion.
\item Achieved 65.75\% top-1 3D grounding success on DOVE-G versus OVIR-3D's 41.5\%, using full queries and a 3D IoU threshold of 0.5.\end{itemize}}
\end{papers}
\cvsection{Teaching experience}
\begin{description}[leftmargin=78pt,labelwidth=70pt,labelsep=8pt,align=left,font=\normalfont,itemsep=3pt,parsep=0pt,topsep=0pt]
\item[Fall 2026] \href{https://kowndinya2000.github.io/cs440-fa26/}{CS440: Introduction to Artificial Intelligence}, \textbf{Course Instructor}, Rutgers University
\item[Spring 2026] CS440: Introduction to Artificial Intelligence, Teaching Assistant, Rutgers University
\item[Fall 2025] CS460: Introduction to Computational Robotics, Teaching Assistant, Rutgers University
\item[Spring 2025] CS440: Introduction to Artificial Intelligence, Teaching Assistant, Rutgers University
\item[Fall 2024] CS440: Introduction to Artificial Intelligence, Teaching Assistant, Rutgers University
\end{description}
\cvsection{Peer review}
\begin{itemize}[leftmargin=13pt,itemsep=2pt,parsep=0pt,topsep=3pt]
\item IEEE International Conference on Robotics and Automation (ICRA) (2026)
\item IEEE/RSJ International Conference on Intelligent Robots and Systems (IROS) (2025, 2026)
\item Conference on Robot Learning (CoRL) (2026)
\item The International Journal of Robotics Research (IJRR) (2026)
\item International Conference on Automated Planning and Scheduling (ICAPS) (2026)
\end{itemize}
\cvsection{Honors \& awards}
\textbf{2022}\enspace Project showcase, IInvenTiv All IITs R\&D Fair. {\small Research project advised by Prof. Kalidas Yeturu, showcased at IIT Delhi.}\par\vspace{4pt}
\textbf{2019}\enspace Gold Medal, Inter IIT Technical Meet 8.0. {\small Ashoka's Tech for Change Challenge at IIT Roorkee; high-preparation event among 20 IITs.}\par\vspace{4pt}
\cvsection{Leadership \& mentoring}
\cventry{Google Developer Student Clubs, IIT Tirupati}{Feb 2019 -- Jul 2020}{Community Lead}{Organized 12 workshops and hackathons for more than 70 students; led 15 mentoring sessions on Flutter and Android.}
\cvsection{Professional service \& activities}
\textbf{2021}\enspace Developed an Android application for IIT Tirupati's convocation.\par\vspace{3pt}
\textbf{2020}\enspace Completed 80 hours of National Service Scheme service; developed the IIT Tirupati NSS website and led its maintenance team.\par\vspace{3pt}
